\section{Pilot Study}

Complex strategic-game simulations are a productive venue for studying emergent LLM-agent phenomena \cite{tang2025dsgbench, wang2025digitalplayer}, including unethical ones such as deception \cite{bakhtin2022diplomacy, park2024deception}. While those environments are more dynamic and less prone to pre-determined outcomes, they are rarely studied for nuclear escalation, ethical reasoning, or the intersections of both.

A recent study, CivBench \cite{chen2026civbench}, enables inspection into open-ended Civilization V gameplay, where nuclear weapon authorization is but one option in late-game. Built on Vox Deorum \cite{chen2025voxdeorum}, CivBench places an LLM strategist into Sid Meier's Civilization V running the Vox Populi community mod, separating strategic reasoning from tactical execution through rule-based modules. At every decision point, Vox Deorum captures input prompt, pre-hoc reasoning tokens and post-hoc rationale (carried into the next turn's prompt as short-term memory). Among dozens of available actions, LLMs can set the \texttt{use-nuke} flavor to express \emph{authorization} to launch nuclear weapon (0 forbids; 100 always launches when tactical conditions are satisfied; default=50).

% I think this is the right place for a slightly more detailed Civ V / Vox Deorum intro than Experiment Design.
% A recent study, CivBench \cite{chen2026civbench}, enables inspection into open-ended gameplay in Sid Meier's Civilization V, a turn-based 4X strategy game where civilizations manage cities, technologies, diplomacy, economy, military forces, and competing victory paths such as science, diplomacy, culture, and domination. In this environment, nuclear weapon authorization is one late-game strategic option among many, embedded in broader pressures such as deterrence, retaliation, conquest, and attempts to prevent another civilization's victory. Built on Vox Deorum \cite{chen2025voxdeorum}, CivBench places an LLM strategist into Civilization V running the Vox Populi community mod, separating strategic reasoning from tactical execution through rule-based modules. At every decision point, Vox Deorum captures the input prompt, pre-hoc reasoning tokens, and post-hoc rationale, which is carried into the next turn's prompt as short-term memory. Among dozens of available actions, LLMs can set the \texttt{use-nuke} flavor to express \emph{authorization} to launch nuclear weapons: 0 forbids use, 100 always authorizes use when tactical conditions are satisfied, and 50 is the default.

% I think we also need one sentence here to preempt reviewer concern about whether Civ V nuclear use is comparable to real-world nuclear use.
% We can say Civ V models nuclear consequences abstractly (city/population damage, fallout, diplomacy, victory pressure), and our claim is about strategic authorization under pressure rather than real-world equivalence.

From parts of CivBench's self-play dataset (1,200 trajectories), our pilot study found emergent episodes where LLMs set \texttt{use-nuke} = 100 without traces of ethical engagement. Nuclear inclination varies by model identity: five models (Claude Sonnet 4.5, Kimi K2.5, GLM 4.7, DeepSeek V3.2, MiniMax-M2.5) pushed \texttt{use-nuke} upward from the default of 50, while only GPT-OSS-120B inclined to move toward restraint. In 72 maximum-escalation decision points, replay with mentioning of real-world impact failed to push \texttt{use-nuke} below the pre-escalation baseline. Models can become more pragmatic, yet ethical engagement was completely absent in post-hoc justification.
